% Immutable cited-input supplement. Target proofs remain embargoed.
% Target: inv:E2.2; tag 01SV.
% Frozen source SHA256 74497bfe327fae25c9cdcb17334b3bf42a9793a9df7c0a4404d97568f2f760f0.
% CompendiumRef is resolved against the cited S6_09 source edition, not silently identified with the current series statement.

% Exact cited input: series:I / inv:D1.2
\begin{proposition}[Vertical and logarithmic tangent bundles]\hypertarget{tag-01SN}{}\label{inv:D1.2}


There is a canonical representable proper flat holomorphic map
$\bar f:X\to\mathcal B$ with $f=\rho\bar f$. If
\[
 \mathcal T=T_X(-\log W),\qquad
 \mathcal V=\ker\bigl(\mathcal T\to
                   \bar f^*T_{\mathcal B}(-p_0)\bigr),
\]
then, for rank-two vector bundles $E,G$ on $\mathcal B$,
\begin{equation}
 \mathcal V=\bar f^*E,\quad E=\OO\oplus\OO(-1),\qquad
 G=R^1\bar f_*\OO_X=\OO\oplus\OO(-1),
\label{pub:I:eq:D-2}
\end{equation}
\begin{equation}
 \hypertarget{tag-01OZ}{}\label{fam:eq:ab-direct-images}
R^0\bar f_*\OO_X=\OO,\qquad R^2\bar f_*\OO_X=\OO(-1).
\end{equation}
and the following is an exact sequence of locally free sheaves on $X$:
\begin{equation}
 \hypertarget{tag-01P0}{}\label{fam:eq:ab-log-tangent}
0\longrightarrow\bar f^*E\longrightarrow\mathcal T
 \longrightarrow\bar f^*\OO(-5)\longrightarrow0.       
\end{equation}
The trivial summands in $E,G$ are generated by the vertical field $\xi$
and the constant integral cocycle $\gamma$, respectively.
\end{proposition}

% Exact cited input: series:I / inv:D1.1
\begin{proposition}[The root base and its explicit cohomology]\hypertarget{tag-01SM}{}\label{inv:D1.1}


Let $\rho:\mathcal B\to B$ be the analytic root curve with orders $3$ and
$4$ at $p_1,p_2$. Write $R_1,R_2$ for its tautological root line bundles.
Its Picard group is infinite cyclic, with generator
$R=R_1\otimes R_2^{-1}$, and
\[
 R_1=R^4,\quad R_2=R^3,\quad \rho^*\OO_B(1)=R^{12},
 \quad K_{\mathcal B}=R^{-7}.
\]
We abbreviate $R^n$ to $\OO_{\mathcal B}(n)$. Choose sections
$x\in H^0(R^3)$ and $y\in H^0(R^4)$ vanishing at $p_2$ and $p_1$,
respectively. For every integer $n$, put
\[
 r_1(n)\equiv n\pmod3,\quad 0\le r_1(n)<3,
 \qquad r_2(n)\equiv3n\pmod4,\quad0\le r_2(n)<4,
\]
\[
 d(n)=\frac{n-4r_1(n)-3r_2(n)}{12}.
\]
Then
\begin{equation}
 \hypertarget{tag-01P3}{}\label{fam:eq:ab-root-degree}
\rho_*\OO_{\mathcal B}(n)=\OO_B(d(n)),\qquad
 H^i(\mathcal B,\OO(n))=H^i(B,\OO(d(n))).                 
\end{equation}
Equivalently, $h^0(\OO(n))$ is the number of pairs
$(a,b)\in\Z_{\ge0}^2$ with $3a+4b=n$, and
$h^1(\OO(n))=h^0(\OO(-n-7))$.
\end{proposition}

% Exact cited input: series:I / fam:st:A.5.2
\begin{theorem}[Tangent cohomology, direct images and miniversality]\hypertarget{tag-032U}{}\label{push:src:E.2.1}
\hypertarget{tag-01RH}{}\label{fam:src:A.5.2}\label{fam:st:A.5.2}\label{inv:D1.3}\label{inv:D2.1}\label{inv:D3.1}\label{cdp:st:E9}
\hypertarget{tag-02YO}{}\label{night:P5.3}\label{stmt:P5.3}
\hypertarget{tag-02YP}{}\label{night:P5.4}\label{stmt:P5.4}

Let $k,l$ be positive integers with $\gcd(3k,4l)=1$, and let $X$
be a sphere member with multiplicities $(3k,4l)$ and the admissible
logarithmic data of \cref{push:section:E.1}. Put $z=kl$ and
$d=3k+4l$. On $\mathfrak B_{k,l}=\mathbb P(3k,4l)$, the vertical
root bundle $E$, the root lift $\bar f$ and its direct images satisfy
\[
 \begin{split}
 E&=\mathcal O\oplus\mathcal O(-z),\\
 R^q\bar f_*\mathcal O_X&=
 \mathcal O,\quad\mathcal O\oplus\mathcal O(-z),\quad\mathcal O(-z),\\
 K_X&=\bar f^*\mathcal O(z-d).
 \end{split}
\]
The tangent cohomology is
$\bigl(h^0,h^1,h^2,h^3\bigr)(X,T_X)=(1,1,1,0)$.
For the transition-arrow convention
$\kappa(\Phi)=(\partial_c\Phi)\Phi^{-1}$, the Kodaira--Spencer
class of $\partial_c$ is $e\otimes[\gamma]$, and the class of
$\partial_u$ is $-(2\pi i u)^{-1}e\otimes[\gamma]$.
Both classes are nonzero. The primary Kuranishi bracket vanishes,
the Kuranishi germ is a smooth reduced disc, and the fixed-data
period family is locally miniversal.

For the standard $(3,4)$ data, $e=\xi$ and the coarse direct
images are
\begin{equation}
 f_*T_X=\mathcal O\oplus\mathcal O(-1),\qquad
 R^1f_*T_X=\mathcal O\oplus\mathcal O(-1)^2,\qquad
 R^2f_*T_X=\mathcal O.
\label{pub:I:eq:D-4}
\end{equation}
\end{theorem}

% Exact cited input: comp:II / inv:E2.1
\begin{proposition}[Credited re-derivation: The structure sheaves of the comb members]\label{inv:E2.1}
Let $X_m$ be any sufficiently small smooth period section obtained by
the $m$-tooth construction of
\citet[Theorem 1.1]{ViaclovskyCombs2026}, with
$m\ge12$, fixed limiting data at the three distinguished values, and
the standard logarithmic-transformation data. Write
$f_m:X_m\to B$. Then
\[
 R^1f_{m*}\OO_{X_m}=\OO(a_m)\oplus\OO(b_m),\quad
 a_m,b_m\ge0,\quad a_m+b_m=2m-1,
\]
\[
 R^2f_{m*}\OO_{X_m}=\OO(2m-1),\qquad
 (h^{0,0},h^{0,1},h^{0,2},h^{0,3})=(1,2m+1,2m,0).      \tag{E.1}\label{invII:eq:E.1}
\]
The individual summands are not specified by $m$ in the construction.
\end{proposition}

